\chapter{Finite Strategic Games and Nash Equilibrium}
\label{sec:p17-finite-strategic-games}

A strategic game specifies each player's available actions and the
payoff to each player at every joint action. Nash equilibrium tests
unilateral deviations while the other players' strategies stay fixed.
This chapter supplies both original existence routes: the countering
correspondence of Nash's 1950 note and the continuous improvement map
of his 1951 paper \cite{Nash1950EquilibriumPoints,Nash1951NonCooperative}.
Their fixed-point premises are proved below through finite simplex
subdivision and a labeling argument.

\begin{hbdefinition}{def:p17-finite-strategic-game}{Normal form, independent mixtures and payoff units}
Let \(N=\{1,\ldots,n\}\), \(n\geq1\). For each player \(i\), let
\(A_i\) be a finite nonempty action set and let
\[
                 u_i:\mathcal A=\prod_{j=1}^n A_j\longrightarrow\mathbb R.
\]
The tuple \((N,(A_i),(u_i))\) is a finite normal-form game.
A pure strategy is an action in \(A_i\). A mixed strategy is a
lottery \(x_i\in\Delta(A_i)\), with the simplex convention of
\cref{def:p17-finite-lottery}; a pure strategy \(a_i\) is identified
with \(\delta_{a_i}\). The mixed-profile space is
\[
                     X=\prod_i\Delta(A_i).
\]
Write \(X_{-i}=\prod_{j\ne i}\Delta(A_j)\) and
\(\mathcal A_{-i}=\prod_{j\ne i}A_j\) for opponent profiles.
Independently drawn actions at \(x=(x_1,\ldots,x_n)\) have joint law
\[
                P_x(a_1,\ldots,a_n)=\prod_i x_i(a_i).
\]
Define
\[
 U_i(x)=\sum_{a\in\mathcal A}u_i(a)\prod_jx_j(a_j),\qquad
 U_i(y_i,x_{-i})=U_i(x_1,\ldots,x_{i-1},y_i,x_{i+1},\ldots,x_n).
\]
The product over no opponents, when \(n=1\), is 1.
Write \(x[j\leftarrow z_j]\) for the profile obtained by replacing
component \(j\) by \(z_j\). The support of \(x_i\) is
\(\{a_i:x_i(a_i)>0\}\).

The payoffs are supplied cardinal consequence utilities.
If they are intended to represent preferences over lotteries of
outcomes, \cref{thm:p17-finite-vnm} supplies the precise preference
axioms. Ordinal rankings of pure outcomes alone do not determine
these mixed comparisons. Probabilistic independence of players'
draws is distinct from the preference independence axiom.
There is no assumption that different players' utility units are
interpersonally comparable or transferable.
\index{strategic game!finite normal form}\index{mixed strategy!independent draws}
\end{hbdefinition}

\begin{hblemma}{thm:p17-mixed-payoffs}{proved=annex:p17.mixed.payoffs}{The unique separately affine expectation}
\(P_x\) is a probability law. Each \(U_i\) is continuous and affine
in each player's mixture separately, and is the unique separately
affine extension of \(u_i\) from the pure profiles. For every \(i\),
\[
           U_i(x)=\sum_{a_i\in A_i}x_i(a_i)U_i(a_i,x_{-i}).
\]
Replacing \(u_i\) by \(\alpha_i u_i+\beta_i\), with
\(\alpha_i>0\), preserves all of player \(i\)'s best responses
and preserves the equilibrium set. Different players may use
different positive affine scales.
\end{hblemma}

\begin{hbdefinition}{def:p17-nash-equilibrium}{Best responses and Nash equilibrium}
For \(x\in X\), define
\[
\begin{aligned}
 v_i(x_{-i})&=\max_{a_i\in A_i}U_i(a_i,x_{-i}),\\
 \operatorname{BR}_i(x_{-i})
    &=\{y_i\in\Delta(A_i):U_i(y_i,x_{-i})=v_i(x_{-i})\},\\
 \mathcal B(x)&=\prod_i\operatorname{BR}_i(x_{-i}).
\end{aligned}
\]
The profile \(x\) is a Nash equilibrium when
\[
                 U_i(x)\geq U_i(y_i,x_{-i})
              \quad\hbox{for every }i\hbox{ and every }y_i\in\Delta(A_i).
\]
Write \(\mathcal E\subseteq X\) for the equilibrium set.
An equilibrium is pure when all its components are certain lotteries.
The opponents are held fixed in every deviation comparison.
\index{Nash equilibrium!best responses}
\end{hbdefinition}

\begin{hblemma}{thm:p17-best-response-geometry}{proved=annex:p17.best.response.geometry}{Support, inequalities and closed-graph countering}
Each best-response set is the nonempty simplex face supported on
the pure actions attaining \(v_i(x_{-i})\).
For \(x\in X\), the following are equivalent:
\[
\begin{aligned}
 x\in\mathcal E
 &\Longleftrightarrow x\in\mathcal B(x)\\
 &\Longleftrightarrow
 U_i(x)\geq U_i(a_i,x_{-i})\quad(i\in N,\ a_i\in A_i)\\
 &\Longleftrightarrow
 x_i(a_i)>0\ \Longrightarrow\
                  U_i(a_i,x_{-i})=v_i(x_{-i})
                      \quad(i\in N,\ a_i\in A_i).
\end{aligned}
\]
The correspondence \(\mathcal B\) has nonempty compact convex
values and a closed graph. The set \(\mathcal E\) is closed
in compact \(X\), and is specified there by finitely many
polynomial weak inequalities. Nonemptiness is proved below.
\end{hblemma}

\section{The fixed-point foundations}

\begin{hbdefinition}{def:p17-simplex-triangulation}{Simplexes, meshes and admissible labels}
For a nonnegative integer \(d\), the standard \(d\)-simplex is
\[
       \Delta_d=\{z\in\mathbb R^{d+1}:z_j\geq0,\ \sum_{j=0}^d z_j=1\}.
\]
A face sets a selection of these coordinates to zero.
A geometric simplex is the convex hull of affinely independent
vertices; its faces are convex hulls of their subsets.
A finite triangulation consists of finitely many geometric simplexes
covering \(\Delta_d\), with every face included and any two members
intersecting in a common face or not at all. Its mesh is the
largest Euclidean diameter of a member.
For a nonempty face \(F\), its barycenter is the average of its
vertices. The barycentric subdivision replaces each simplex by
the simplexes spanned by barycenters of chains of its nonempty faces.

A Sperner labeling assigns each triangulation vertex a label in
\(\{0,\ldots,d\}\), with label \(j\) allowed only when the
vertex's \(j\)-th coordinate is positive. A top-dimensional
simplex is fully labeled when its vertices have all \(d+1\) labels.
\end{hbdefinition}

\begin{hblemma}{thm:p17-sperner-foundation}{proved=annex:p17.sperner.foundation}{Subdivision, shrinking mesh and the complete parity lemma}
Barycentric subdivision is a finite triangulation and respects
all boundary faces. For \(d\geq1\), its mesh is at most
\(d/(d+1)\) times the preceding mesh. Iterating gives
triangulations with mesh tending to zero.
Every Sperner labeling has an odd number of fully labeled
top-dimensional simplexes, and therefore at least one.
For \(d=0\), the sole vertex supplies the complete assertion.
\end{hblemma}

\begin{hbtheorem}{thm:p17-brouwer-foundation}{proved=annex:p17.brouwer.foundation}{Brouwer on every nonempty compact convex finite-dimensional domain}
Every continuous map \(f:\Delta_d\to\Delta_d\) has a fixed point.
More generally, if \(K\) is a nonempty compact convex subset of
\(\mathbb R^d\), every continuous \(f:K\to K\) has a fixed point.
This includes a singleton and sets of lower affine dimension.
The proof gives the simplex labeling argument and the continuous
nearest-point projection needed to pass to general \(K\).
\end{hbtheorem}

\begin{hbtheorem}{thm:p17-kakutani-foundation}{proved=annex:p17.kakutani.foundation}{The complete closed-graph correspondence theorem needed by Nash}
Let \(K\subseteq\mathbb R^d\) be nonempty, compact and convex.
Suppose \(\mathcal C(x)\subseteq K\) is nonempty, closed and
convex for each \(x\in K\), and its graph
\[
                   \{(x,y)\in K\times K:y\in\mathcal C(x)\}
\]
is closed. Then some \(x\in K\) satisfies \(x\in\mathcal C(x)\).
The proof constructs continuous approximate selections from finite
nets, uses the preceding Brouwer theorem, and proves the limiting
membership using closed graph and convexity.
\end{hbtheorem}

\section{The two original existence constructions}

\begin{hbtheorem}{thm:p17-nash-1950}{proved=annex:p17.nash.1950}{Nash's countering-correspondence proof}
Every finite normal-form game has an independent mixed Nash
equilibrium. Nash's 1950 countering correspondence is precisely
\(\mathcal B\): a profile \(y\) counters \(x\) when every \(y_i\)
is a best response to \(x_{-i}\). Its fixed points are exactly
the equilibrium profiles. All the premises of
\cref{thm:p17-kakutani-foundation} hold on \(X\)
\cite[pp.~48--49]{Nash1950EquilibriumPoints}.
\end{hbtheorem}

\begin{hbdefinition}{def:p17-nash-improvement-map}{Nash's 1951 continuous map}
For \(i\in N\), \(a_i\in A_i\), and \(x\in X\), define
\[
\begin{aligned}
 g_{i,a_i}(x)&=\max\{0,U_i(a_i,x_{-i})-U_i(x)\},\\
 G_i(x)&=\sum_{a_i\in A_i}g_{i,a_i}(x),\\
 T_i(x)(a_i)&=\frac{x_i(a_i)+g_{i,a_i}(x)}{1+G_i(x)},\qquad
                          T(x)=(T_i(x))_{i\in N}.
\end{aligned}
\]
This is the original normalized positive-gain construction,
with the source's \(\varphi\) written as \(g\)
\cite[p.~288, theorem~1]{Nash1951NonCooperative}.
\end{hbdefinition}

\begin{hbtheorem}{thm:p17-nash-1951}{proved=annex:p17.nash.1951}{Nash's continuous-map proof and its exact fixed points}
The map \(T:X\to X\) is continuous, and
\[
                              T(x)=x\quad\Longleftrightarrow\quad x\in\mathcal E.
\]
Therefore \cref{thm:p17-brouwer-foundation} gives the existence
theorem independently of the countering-correspondence application.
The claim is about the existence of a fixed point; it does not
assert convergence of repeated application of \(T\).
\end{hbtheorem}

\begin{hbdefinition}{def:p17-game-symmetry}{The source's full symmetry group}
A game automorphism consists of a permutation \(\sigma\) of
players and bijections \(h_i:A_i\to A_{\sigma(i)}\), such that
for every pure profile \(a\),
\[
                     u_{\sigma(i)}(h(a))=u_i(a)\quad(i\in N),
\]
where the \(\sigma(i)\)-th component of \(h(a)\) is \(h_i(a_i)\).
It acts linearly on mixtures by permuting their coordinates.
A profile is symmetric when every game automorphism fixes it.
This definition does not identify strategies of different players
unless an actual automorphism supplies that identification.
\end{hbdefinition}

\begin{hbtheorem}{thm:p17-nash-symmetry}{proved=annex:p17.nash.symmetry}{An equilibrium fixed by every game symmetry}
Every finite game has an equilibrium invariant under all its
automorphisms. The invariant-profile set is a nonempty compact
convex set and is preserved by \(T\)
\cite[p.~289, theorem~2]{Nash1951NonCooperative}.
\end{hbtheorem}

\section{The original solution-set notions}

\begin{hbdefinition}{def:p17-nash-solution-sets}{Interchangeability, solutions and subsolutions}
A nonempty \(C\subseteq\mathcal E\) is interchangeable when
\[
          x,y\in C\quad\Longrightarrow\quad(y_i,x_{-i})\in C
                                                    \quad(i\in N).
\]
The game is solvable in Nash's 1951 terminology when its entire
\(\mathcal E\) is interchangeable; its solution is then
\(\mathcal E\). A subsolution is an inclusion-maximal interchangeable
subset of \(\mathcal E\). Write \(C_i=\{x_i:x\in C\}\) for its factors.
A solution is strong when
\[
 x\in\mathcal E,\quad U_i(y_i,x_{-i})=U_i(x)
                      \quad\Longrightarrow\quad(y_i,x_{-i})\in\mathcal E.
\]
These are the source's solution-set notions, distinct from the
existence of an individual Nash equilibrium
\cite[pp.~290--291]{Nash1951NonCooperative}.
\end{hbdefinition}

\begin{hbtheorem}{thm:p17-nash-subsolutions}{proved=annex:p17.nash.subsolutions}{Existence, products, compact factors and values}
Every equilibrium belongs to a subsolution. Every subsolution
satisfies
\[
                              C=\prod_i C_i,
\]
and every \(C_i\) is nonempty, compact and convex.
There is a unique subsolution exactly when the game is solvable;
it then equals \(\mathcal E\).
The subsolution existence assertion uses Zorn's lemma.

The continuous payoffs attain maxima and minima on
\(\mathcal E\) and on each subsolution. Nash's upper and lower
equilibrium values are
\[
                   v_i^+=\max_{x\in\mathcal E}U_i(x),\qquad
                   v_i^-=\min_{x\in\mathcal E}U_i(x).
\]
A common value exists when these are equal. Uniqueness of an
equilibrium implies equality, but is not required for it.
This gives the full product and closed-convex factor conclusions
of source Theorems~3--4.
\end{hbtheorem}

\begin{hbtheorem}{thm:p17-nash-polyhedral}{proved=annex:p17.nash.polyhedral}{Finite half-space descriptions for solvable games}
If the game is solvable, fix any \(t\in\mathcal E\).
Its \(j\)-th equilibrium factor is exactly
\[
 \left\{z_j\in\Delta(A_j):
 U_i(t[j\leftarrow z_j])-
 U_i(t[j\leftarrow z_j][i\leftarrow a_i])\geq0
                 \quad(i\in N,\ a_i\in A_i)\right\}.
\]
The further replacement includes \(i=j\).
Every displayed constraint is affine in \(z_j\).
Thus each factor is a bounded polyhedron, and is the convex hull
of finitely many vertices
\cite[p.~292, theorem~5]{Nash1951NonCooperative}.
\end{hbtheorem}

\begin{hbtheorem}{thm:p17-finite-minimax}{proved=annex:p17.finite.minimax}{The exact two-player zero-sum consequence}
For a real \(r\times s\) matrix \(M\), \(r,s\geq1\), give
the players payoffs \(x^{\mathsf T}My\) and
\(-x^{\mathsf T}My\), where \(x\in\Delta(\{1,\ldots,r\})\)
and \(y\in\Delta(\{1,\ldots,s\})\). Then
\[
 \max_x\min_y x^{\mathsf T}My
   =\min_y\max_x x^{\mathsf T}My=v.
\]
Both extrema are attained. Let
\[
\begin{aligned}
 X^*&=\{x:x^{\mathsf T}Me_b\geq v\text{ for every column }b\},\\
 Y^*&=\{y:e_a^{\mathsf T}My\leq v\text{ for every row }a\}.
\end{aligned}
\]
Here \(e_a,e_b\) are the corresponding pure unit vectors,
and \(x,y\) remain in their stated probability simplexes.
Then \(\mathcal E=X^*\times Y^*\), and every equilibrium gives
the same payoff pair \((v,-v)\). Conversely the displayed
minimax equality with attainment gives an equilibrium.
Hence the existence and minimax conclusions are equivalent in
this two-player zero-sum domain, as stated in the 1950 note.
Such games are solvable. Strong solvability implies a pure saddle
point; the existence of a pure saddle alone need not imply strong
solvability.
\end{hbtheorem}

\section{Complete finite examples and deviation boundaries}

\begin{hbdefinition}{def:p17-nash-six-games}{Nash's six two-by-two examples}
Rows are \(a,b\), columns are \(\alpha,\beta\), and each entry lists
the row and column payoff in that order. These are the six source
games on pp.~291--292:
\[
\begin{array}{c|cc}
1&\alpha&\beta\\ \hline
a&(5,-3)&(-4,4)\\ b&(-5,5)&(3,-4)
\end{array}
\qquad
\begin{array}{c|cc}
2&\alpha&\beta\\ \hline
a&(1,1)&(-10,10)\\ b&(10,-10)&(-1,-1)
\end{array}
\]
\[
\begin{array}{c|cc}
3&\alpha&\beta\\ \hline
a&(1,1)&(-10,-10)\\ b&(-10,-10)&(1,1)
\end{array}
\qquad
\begin{array}{c|cc}
4&\alpha&\beta\\ \hline
a&(1,1)&(0,1)\\ b&(1,0)&(0,0)
\end{array}
\]
\[
\begin{array}{c|cc}
5&\alpha&\beta\\ \hline
a&(1,2)&(-1,-4)\\ b&(-4,-1)&(2,1)
\end{array}
\qquad
\begin{array}{c|cc}
6&\alpha&\beta\\ \hline
a&(1,1)&(0,0)\\ b&(0,0)&(0,0).
\end{array}
\]
Write \(p=x_1(a)\) and \(q=x_2(\alpha)\), so a profile is
specified by \((p,q)\in[0,1]^2\).
\end{hbdefinition}

\begin{hbtheorem}{thm:p17-nash-six-games}{proved=annex:p17.nash.six.games}{Every equilibrium and payoff in the six source games}
The equilibrium sets and corresponding payoffs are:
\[
\begin{array}{c|c|c}
\text{game}&(p,q)&(U_1,U_2)\\ \hline
1&(9/16,7/17)&(-5/17,1/2)\\
2&(0,0)&(-1,-1)\\
3&(1,1),(0,0),(1/2,1/2)&(1,1),(1,1),(-9/2,-9/2)\\
4&\text{every }(p,q)&(q,p)\\
5&(1,1),(0,0),(1/4,3/8)&(1,2),(2,1),(-1/4,-1/4)\\
6&(1,1),(0,0)&(1,1),(0,0).
\end{array}
\]
Game 1 has no pure equilibrium. Games 1, 2 and 4 are solvable;
games 2 and 4 are strongly solvable, while game 1 is not.
Games 3, 5 and 6 fail interchangeability.
In game 3 the uniform strategy both maximizes its minimum
payoff and minimizes its maximum payoff against pure opponents,
with value \(-9/2\); the players' actual payoffs remain identical.
In game 6, at a perturbed profile
\((\varepsilon,\varepsilon)\), \(0<\varepsilon<1\), each player
can improve by \(\varepsilon(1-\varepsilon)\) by choosing its
first action certainly. This is a precise perturbation calculation,
without an unstated adjustment dynamics.
\end{hbtheorem}

\begin{hbdefinition}{def:p17-strategic-dominance}{Strict domination of a player's strategy}
A strategy \(y_i\) strictly dominates \(x_i\) when
\[
               U_i(y_i,z_{-i})>U_i(x_i,z_{-i})
                          \quad\hbox{for every }z_{-i}\in X_{-i}.
\]
This compares two strategies of one player against all opponent
profiles. It is not the coalitional domination relation on
imputations used for NM solutions.
\end{hbdefinition}

\begin{hblemma}{thm:p17-strategic-dominance}{proved=annex:p17.strategic.dominance}{Finite domination witnesses and safe exclusions}
It suffices to check the strict domination inequality at all
pure opponent profiles. A strictly dominated strategy cannot
be a component of a Nash equilibrium.
If \(y_i\) strictly dominates \(x_i\), and \(z_i(a)>0\)
whenever \(x_i(a)>y_i(a)\), then for sufficiently small
\(\rho>0\),
\[
                       z_i'=z_i+\rho(y_i-x_i)
\]
is a mixed strategy and strictly dominates \(z_i\).
One may take any positive \(\rho\) at most
\[
       \min_{a:x_i(a)>y_i(a)}
                      \frac{z_i(a)}{x_i(a)-y_i(a)}.
\]
In particular the undominated strategies form a union of faces
of the strategy simplex.
\end{hblemma}

\begin{hbdefinition}{def:p17-correlated-recommendation}{Joint recommendations and obedience}
A recommendation law is any \(\mu\in\Delta(\mathcal A)\).
A device draws \(a\) from \(\mu\) and privately recommends
\(a_i\) to player \(i\). The obedience inequalities are
\[
 \sum_{a_{-i}}\mu(a_i,a_{-i})
       [u_i(a_i,a_{-i})-u_i(b_i,a_{-i})]\geq0
                          \quad(i,\ a_i,b_i\in A_i).
\]
These inequalities define a correlated equilibrium for this
finite recommendation model. They are unconditional weighted
inequalities, so no conditioning on a zero-probability
recommendation is performed.
\end{hbdefinition}

\begin{hbtheorem}{thm:p17-product-and-correlation}{proved=annex:p17.product.and.correlation}{The exact independence boundary}
The obedience inequalities are equivalent to absence of gain
from any randomized deviation depending only on one's own
recommendation. A product law \(P_x\) obeys them if and only if
\(x\) is a Nash equilibrium.
The law putting probability \(1/2\) on each diagonal outcome
of game 3 is obedient and gives payoff \((1,1)\). It is not
a product law. Its marginal product is the mixed Nash
equilibrium with payoff \((-9/2,-9/2)\).
Thus equal individual marginals do not determine the incentives
or payoffs of a joint recommendation scheme.
\end{hbtheorem}

\begin{hbtheorem}{thm:p17-strategic-welfare-boundaries}{proved=annex:p17.strategic.welfare.boundaries}{Equilibrium, Pareto comparisons, welfare and cooperative solution objects}
Nash equilibrium need not be Pareto efficient: game 2 has
equilibrium payoff \((-1,-1)\), strictly improved for both
players at \((a,\alpha)\).
Here Pareto domination means a weak payoff improvement for
every player and a strict improvement for at least one.

Player-specific positive affine rescaling preserves both the
equilibrium set and Pareto comparisons, but need not preserve
rankings by the sum of payoffs. In game 5, multiplying only
player 1's utilities by 2 makes \((b,\beta)\) have the greater
sum among the two pure equilibria; multiplying only player 2's
utilities by 2 makes \((a,\alpha)\) have the greater sum.
An aggregate-welfare comparison therefore needs additional
interpersonal scales or weights.

A Nash equilibrium is a profile tested against unilateral
strategy changes. An NM solution is a set of efficient,
individually rational imputations tested against feasible
coalitional domination, containing no internally dominating pair,
and required to dominate all outsiders.
The exact finite-game existence theorem above and the
ten-player nonexistence theorem
\cref{thm:p18-lucas-positive-cover} concern these different
objects and are compatible. Nash's separate interchangeability
notion of a \emph{solution} must also retain its stated definition.
\end{hbtheorem}

\begin{hbdefinition}{def:p17-finite-nash-approximation}{A finite procedure for a stated accuracy}
For rational payoff data and rational \(\varepsilon>0\), let
\[
 M=\max_{i,a}|u_i(a)|,\qquad L=\sum_i(|A_i|-1).
\]
If \(ML=0\), choose any pure profile. Otherwise choose an integer
\(k\geq4ML/\varepsilon\). Enumerate all profiles whose
probabilities are nonnegative multiples of \(1/k\).
For each, calculate the regret
\[
                    r(x)=\max_{i,a_i}[U_i(a_i,x_{-i})-U_i(x)].
\]
Return the first profile with \(r(x)\leq\varepsilon\).
Such a profile is called an \(\varepsilon\)-Nash equilibrium.
\end{hbdefinition}

\begin{hbtheorem}{thm:p17-finite-nash-approximation}{proved=annex:p17.finite.nash.approximation}{Termination and the exact approximation guarantee}
The procedure terminates and returns an \(\varepsilon\)-equilibrium.
If \(ML=0\), it returns a pure equilibrium immediately.
Otherwise it examines at most
\[
                         \prod_i {k+|A_i|-1\choose |A_i|-1}
\]
profiles, using exact rational arithmetic. This is a finite
accuracy guarantee; no efficient bound or exact-equilibrium
computation is inferred from it.
\index{Nash equilibrium!finite approximation}
\end{hbtheorem}

The source's three-person poker analysis uses behavioral parameters
at different information situations and refers to a separate
Nash--Shapley treatment. Its full behavioral/normal-form bridge,
elimination argument and payoff calculation form a further
presentation. The proposed negotiation models likewise require
their own game specification and, when infinitely many strategies
are admitted, their own existence hypotheses.

For the beliefs that can support a strategic choice, use
\cref{sec:p17-epistemic-game-theory}. It proves the finite
common-belief/elimination correspondence and gives exact cases
in which rationalizable choices exceed the choices supported
by Nash equilibrium.
